\begin{abcliste}
\abc Electron and proton carry the elementary charge $e$, with opposite signs. With $r$ the distance, they attract each other with
\begin{align}
 F_\mathrm{el} &= \FelF \\
   &= \kC\times\frac{\left(\qe\right)^2}{\left(\rH\right)^2} \\
   &= \Fel \approx \result{\FelP{-9}{2}}
\end{align}

\abc The gravitational force is $F_\mathrm{G} = G\,m_e\,m_p/r^2$. Both forces fall off with $1/r^2$, so the distance cancels in the ratio:
\begin{align}
 \frac{F_\mathrm{el}}{F_\mathrm{G}} &= \ratioF \\
   &= \frac{\kC\times\left(\qe\right)^2}{\G\times\me\times\mpr} \\
   &= \ratio \approx \result{\ratioP{0}{2}}
\end{align}
In atoms, gravity plays no part at all: the electric force holds them together.
\end{abcliste}
